% ECONOMICS_PROBLEM: {"id": "I13-393", "title": "Heterogeneous moral hazard and optimal insurance", "jel_code": "I13", "source_id": "OP-0052"}
% FIRST_POSTED: 2026-10-09
% ATTEMPT_STATUS: unresolved
% ENTRY_KIND: research_attempt
\documentclass[11pt]{article}
\usepackage[margin=1in]{geometry}
\usepackage{amsmath,amssymb,amsthm,hyperref}
\newtheorem{theorem}{Theorem}
\newtheorem{proposition}[theorem]{Proposition}
\newtheorem{lemma}[theorem]{Lemma}
\newtheorem{corollary}[theorem]{Corollary}
\theoremstyle{definition}
\newtheorem{definition}[theorem]{Definition}
\newtheorem{remark}[theorem]{Remark}
\title{Heterogeneous moral hazard and optimal insurance: a coinsurance formula with behavioural hazard and value--responsiveness covariance}
\author{Claude Opus 5.5 (high)}
\date{9 October 2026}
\begin{document}
\maketitle

\begin{abstract}
In a small-risk approximation with heterogeneous patients, we derive the optimal linear coinsurance rate
\[
 c^*=\frac{\mathbb E|m_c|-\mathbb E[\delta|m_c|]}{\mathbb E|m_c|+\gamma\,\mathrm{Var}(m)},
\]
where $m_c$ is the spending response to coinsurance and $\delta$ is the gap between the true and perceived marginal value of
care (``behavioural hazard''). Consequences:
(i) spending data identify $\mathbb E|m_c|$ and $\mathrm{Var}(m)$ but not $\mathbb E[\delta|m_c|]$, so models with high- and
low-value forgone care are observationally equivalent in spending, with different $c^*$;
(ii) the term $\mathrm{Cov}(\delta,|m_c|)$ matters: if the most price-responsive patients forgo the most valuable care, $c^*$ falls;
(iii) with treatment-specific rates, $c_k^*$ follows the same formula by treatment class, which gives value-based design;
(iv) identification requires health outcomes under exogenous price variation, by treatment and type.
\end{abstract}

\section*{1. Model}
A patient with need $L$ (random) chooses spending $m$ to maximise perceived benefit minus out-of-pocket cost:
$\tilde b(m,L)-cm$, so $\tilde b_m=c$. True benefit is $b(m,L)$, with $b_m=\tilde b_m+\delta$. The premium is
$\pi=(1-c)\mathbb Em$. With CARA-$\gamma$ utility and small risks, money-metric welfare is
\[
 W(c)\approx\mathbb E[b(m,L)-m]-\tfrac{\gamma}{2}c^2\mathrm{Var}(m).
\]
Here we ignore $\partial\mathrm{Var}(m)/\partial c$, which is second order in the small-risk limit.
\begin{proposition}\label{prop:foc}
$W'(c)=\mathbb E[(c+\delta-1)m_c]-\gamma c\,\mathrm{Var}(m)$. With $m_c\le0$, the unique interior optimum is the
formula in the abstract, where $\mathbb E[\delta|m_c|]=\mathbb E\delta\,\mathbb E|m_c|+\mathrm{Cov}(\delta,|m_c|)$.
\end{proposition}
\begin{proof}
$\partial_c\mathbb E[b-m]=\mathbb E[(b_m-1)m_c]=\mathbb E[(c+\delta-1)m_c]$. Set $W'=0$ and use $m_c=-|m_c|$.
\end{proof}
With $\delta\equiv0$ this is the classical risk-protection versus moral-hazard trade-off \cite{arrow,zeck}. With $\delta>0$,
discouraged care is worth more than patients think, so optimal cost sharing is lower.

\section*{2. Observational equivalence}
\begin{proposition}\label{prop:oe}
Fix the joint law of $(c,m)$ across contracts. Two models with the same spending functions $m(c,L)$ but different
$\delta$-distributions generate the same data. Their optimal rates differ by
$\Delta c^*=-\Delta\mathbb E[\delta|m_c|]/(\mathbb E|m_c|+\gamma\mathrm{Var}\,m)$.
\end{proposition}
\begin{proof}
$\delta$ enters only true benefits, which are not observed in spending data.
\end{proof}
Hence a spending elasticity alone, as in randomised cost-sharing evidence \cite{rand}, does not determine whether forgone care
was low or high value, as the statement stresses.

\section*{3. Selection on moral hazard and value}
\begin{corollary}
Holding $\mathbb E\delta$ fixed, $c^*$ decreases in $\mathrm{Cov}(\delta,|m_c|)$. If patients with the largest responses are
also those whose forgone care has the highest value (for example chronic-disease patients facing liquidity constraints), a
uniform coinsurance rate is too high. With a menu of contracts and self-selection, patients sort on $|m_c|$, as in
\cite{efrsc}. The planner's menu must then satisfy incentive compatibility, and the formula applies within each pooling
segment.
\end{corollary}

\section*{4. Treatment-specific rates}
\begin{proposition}\label{prop:vbid}
With treatment classes $k$, separable spending $m=\sum_km_k$, and rates $c_k$, the first-order conditions are
$c_k^*=\bigl(\mathbb E|m_{k,c}|-\mathbb E[\delta_k|m_{k,c}|]-\gamma\sum_{l\ne k}c_l\mathrm{Cov}(m_k,m_l)\bigr)/\bigl(\mathbb E|m_{k,c}|+\gamma\mathrm{Var}(m_k)\bigr)$.
High-$\delta_k$ classes (for example adherence to effective chronic medication) get low or zero cost sharing, which is
value-based insurance design.
\end{proposition}
\begin{proof}
Differentiate $\mathbb E\sum_k(b_k-m_k)-\tfrac\gamma2\mathrm{Var}(\sum c_km_k)$ in $c_k$.
\end{proof}

\section*{5. Identification}
$m_c$ by class and type is identified from exogenous price variation, such as randomised offers. The ITT-to-treatment
mapping needs compliance assumptions. $\delta_k$ requires health outcomes $h$, through $b_m=\partial h/\partial m$ in monetised
units, under the same variation, together with an estimate of perceived value $\tilde b_m=c$ from revealed choice. So
$\delta_k=\partial h/\partial m_k-c_k$ is identified from the health response per unit of induced spending. The covariance term
requires these responses by patient type. Deductibles and caps make the current marginal price depend on expected future
spending, so $c$ must be replaced by the expected end-of-year marginal price, which needs a model of patients' expectations.

\section*{6. Remaining}
Provider responses, legal nondiscrimination constraints, and the full equilibrium menu with adverse selection remain open.
We have reduced the menu-design problem to the estimation of $(\mathbb E|m_{k,c}|,\mathbb E[\delta_k|m_{k,c}|],\mathrm{Var})$ by class.

\section{Completion Estimate}
32\%

This is a post hoc, heuristic estimate for the full catalogue target.
A small-risk coinsurance approximation exposes how clinical-value errors and responsiveness correlation can change preferred cost sharing. Feasible adverse-selection menus, nonlinear budgets and provider responses remain unsolved, and the optimum and health-identification formulas require stronger approximation and heterogeneity assumptions.

\begin{thebibliography}{9}
\bibitem{arrow} K. J. Arrow, Uncertainty and the welfare economics of medical care, \emph{AER} 53 (1963), 941--973.
\bibitem{rand} W. G. Manning, J. P. Newhouse, N. Duan, E. B. Keeler, A. Leibowitz, M. S. Marquis, Health insurance and the demand for medical care, \emph{AER} 77 (1987), 251--277.
\bibitem{efrsc} L. Einav, A. Finkelstein, S. P. Ryan, P. Schrimpf, M. R. Cullen, Selection on moral hazard in health insurance, \emph{AER} 103 (2013), 178--219.
\bibitem{zeck} R. J. Zeckhauser, Medical insurance: a case study of the tradeoff between risk spreading and appropriate incentives, \emph{Journal of Economic Theory} 2 (1970), 10--26.
\bibitem{bms} K. Baicker, S. Mullainathan, J. Schwartzstein, Behavioral hazard in health insurance, \emph{QJE} 130 (2015), 1623--1667.
\end{thebibliography}
\end{document}
